\section*{Problema 3}
Usando solo las cuatro propiedades del producto escalar, verifica en detalle las identidades dadas en el texto para $(A + B)^2$ y $(A - B)^2$.

\section*{Problema 4}
¿Cuál de los siguientes pares de vectores son perpendiculares?
\begin{itemize}
    \item (a) $(1, -1, 1)$ y $(2, 1, 5)$
    \item (b) $(1, -1, 1)$ y $(2, 3, 1)$
    \item (c) $(-5, 2, 7)$ y $(3, -1, 2)$
    \item (d) $(\pi, 2, 1)$ y $(2, -\pi, 0)$
\end{itemize}

\section*{Problema 5}
Considera funciones continuas en el intervalo $[-1, 1]$. Define el producto escalar de dos funciones $f, g$ como:
\[ \int_{-1}^{1} f(x)g(x) \, dx. \]

Denotamos esta integral también como $\langle f, g \rangle$. Verifica que las cuatro reglas para un producto escalar se cumplen, en otras palabras, demuestra que:
\begin{itemize}
    \item SP 1. $\langle f, g \rangle = \langle g, f \rangle$.
    \item SP 2. $\langle f, g + h \rangle = \langle f, g \rangle + \langle f, h \rangle$.
    \item SP 3. $\langle cf, g \rangle = c \langle f, g \rangle$.
    \item SP 4. Si $f = 0$, entonces $\langle f, f \rangle = 0$ y si $f \neq 0$, entonces $\langle f, f \rangle > 0$.
\end{itemize}

\section*{Problema 6}
Si $f(x) = x$ y $g(x) = x^2$, ¿cuáles son $\langle f, f \rangle$, $\langle g, g \rangle$ y $\langle f, g \rangle$?

\section*{Problema 7}
Considera funciones continuas en el intervalo $[-\pi, \pi]$. Define un producto escalar similar al anterior para este intervalo. Demuestra que las funciones $\sin nx$ y $\cos mx$ son ortogonales para este producto escalar ($m, n$ siendo enteros).

\section*{Problema 8}
Sea $A$ un vector perpendicular a cada vector $X$. Demuestra que $A = 0$.